\documentclass{article}
\usepackage[utf8]{inputenc}
\usepackage[spanish]{babel}
\usepackage{listings}
\usepackage{graphicx}
\graphicspath{ {images/} }
\usepackage{cite}

\begin{document}

\begin{titlepage}
    \begin{center}
        \vspace*{1cm}
            
        \Huge
        \textbf{Primer Informe}
            
        \vspace{0.5cm}
        \LARGE
        Calistenia
            
        \vspace{1.5cm}
            
        \textbf{Carlos García Guerra}
            
        \vfill
            
        \vspace{0.8cm}
            
        \Large
        Despartamento de Ingeniería Electrónica y Telecomunicaciones\\
        Universidad de Antioquia\\
        Medellín\\
        Marzo de 2021
            
    \end{center}
\end{titlepage}


\newpage
\section{Introducción al desafío}
Este desafío consiste en seguir una serie de instrucciones o pasos para conseguir llevar dos tarjetas de una posición A, a una posición B. El reto está en que yo planteo una solución, establezco unas instrucciones a seguir, tratando de ser lo más preciso, evitando ambigüedades. Sin embargo lo interesante de todo esto es que los pasos por más precisos que sean, fueron escritos bajo mi criterio, estimando que se llevarán al pie de la letra según mi propia interpretación de dichos pasos, pero se vuelve complejo en este punto ya que cada persona los puede interpretar de una forma distinta.
\section {Solución desafío: Llevar unos objetos de una posición A a una posición B.}
1. Partiendo de la posición inicial, una hoja de papel puesta sobre una superficie plana (En su defecto una mesa) y dos tarjetas del mismo grosor y tamaño sobre la misma.\\

2. Con la mano que tengas más habilidad, (solo con una mano) agarra las dos tarjetas y ponlas una encima de la otra, en forma horizontal, dejándolas lo más alineadas posible.\\

3.  De ese mismo modo tome ambas tarjetas y pase a ubicarlas en la mitad de la hoja de papel, dejándolas nuevamente alienadas, una encima de otra. (Si ya las había puesto en el centro de la hoja en el anterior paso, salte al paso 4).\\

4.  Recuerde que todos lo pasos son con una sola mano. Ahora pondrá ambas tarjetas de manera vertical a la superficie solo siguiendo estas indicaciones: ponga su mano en forma de garra y ubiquela sobre las tarjetas del tal forma en que su dedo pulgar quede posicionado al lado del extremo más largo y el resto de dedos al otro costado de la tarjeta. Ahora colóquelas de forma vertical formando un angulo de 90° con la superficie.\\

5. En este punto, manteniendo la posición, su dedo indice ubiquelo por encima de las tarjetas y haga presión con ese dedo hacia abajo, mientras que con el dedo pulgar y anular, trata de separar las tarjetas una de otra.\\

6. Una vez separadas, trate de estabilizar ambas tarjetas en especie de pirámide y usar los dedos mencionados en el paso anterior para ayudarse a hacerlas sostener sobre la superficie con paciencia.


\bibliographystyle{IEEEtran}

\end{document}
